\documentclass[12pt]{article}
\begin{document}


\section{Class Questions}

$\in$: `is a' or `is an element of'
$$a \in \{a, b, c\}$$

$\{a\} \in \{a, b, c\}$?
There are many canine pets, let's call the collection of them dogs.
My pets name is Maeby, she is a dog, that is, Maeby is an element of the set of dogs.
Is the collection that only includes Maeby a dog?

Go through example 1.3 in Book of Proof.



\subsection{Breakout Sessions}
Which of the following are correctly constructed?
$$a \in \{\{a, b\}, c\}$$
$$c \in \{a,b,\{c\}, c\}$$
$$\{a,b\} \in \{\{a, b\}, c\}$$
$$\{b\} \in \{\{a, b\}, c\}$$
$$[4, 3) \in \mathbf{R}$$

Sketching Functions

$$f(x)=x^3-2x+8$$

$$f(x)=3x^2-log(x)+2$$

$$f(x)=1-2*log(x/2)$$

$$f(x)=\frac{1}{3x^2}+2$$

\pagebreak
\section{Set Equality}

We introduced in lecture the concept of a set, and the constructor notation
e.g. $A=\{x: x \in [0,1]\}$

The following sounds pedantic:\\
\begin{itemize}
\item A set $A$ is = to a set $B$ if they have exactly the same elements, $x\in A$ iff $x\in B$.\\
\item A set $A$ is subset of $B$ if every element of $A$ is in $B$.\\
\end{itemize}
But from here it follows that if A set $A$ is subset of $B$ and $B$ is subset of $A$, then $A=B$ and vice versa.

How do we know?  
Start from the left, use the definitions:
Pick an arbitrary element of A, by the definition of subset, it is in B. So  $x\in A$ $\rightarrow$ $x\in B$
Pick an arbitrary element of B, by the definition of subset, it is in A.  So  $x\in B$ $\rightarrow$ $x\in A$.
We are all set.

\subsection{Example 1: Distributive property}
A non-empty collection of sets is called a "class" or "family"; where $A\in\mathcal{A}$ are sets.

The union of all member of a class is denoted $\cup\mathcal{A}$
$$\cup\{A : A\in \mathcal{A}\}$$


Let $A$ be a set and $\mathcal{B}$ be a class of sets:
Prove that 
$$A\cap \cup\mathcal{B}=\bigcup\{A\cap B: B\in \mathcal{B}\}$$
$$Z=Q$$
We are asked to prove that two sets are equal, they are if the first is a subset the second and the second is a subset of the first.

Part 1 $A\cap \cup\mathcal{B}\subseteq \bigcup\{A\cap B: B\in \mathcal{B}\}$

Pick $x \in Z$, we want to show that it is in $Q$

$x \in Z$, so $x\in A$ and $x\in B_i$ call that $B_1$ 
$$(A\cap B_1)$$
so an element of a set is also the element of that set union some other set:
$$x\in(A\cap B_1)\cup(A\cap B_2)\ldots$$
$$\bigcup\{A\cap B: B\in \mathcal{B}\}$$



Part 2: $A\cap \cup\mathcal{B}\supseteq \bigcup\{A\cap B: B\in \mathcal{B}\}$

Pick $x \in Q$, we want to show that it is in $Z$.  To do that we need to show that x is in A and x is in $\cup\mathcal{B}$.

x in Q means that it is in $(A\cap B_i)$ for some B.  
This means that x is in A.
x is also in $B_i$, so x is in $\cup\mathcal{B}$

All done!
\pagebreak
\section{Functions}
Recall we refer to all the combinations of two sets $X$ and $Y$ with $X\times Y$.

A function $f$ that maps $X$ into $Y$ $f: X\rightarrow Y$
is a relation $f\subseteq X\times Y$ such that
\begin{itemize}
\item for every $x\in X$, there exists a $y \in Y$ such that $f(x)=y $.
\item For every $y,\ z \in Y$ with $f(x)=y $ and $f(x)=z $, we have $y=z$
\end{itemize}
X is called the domain of $f$; Y is the codomain of $f$; the range of $f$ is
$$f(X):=\{y\in Y: f(x)=y \textrm{ for some }x \in X\}$$

General Rule:
If X and Y are a nonempty sets, $f: X\rightarrow Y$, then for any $\mathcal{A}\subseteq 2^X$

$$f(\cup\mathcal{A})= \bigcup \{f(A): A\in \mathcal{A}\}$$

E.g.

$f=x^2$

$X=\{1, 2, 3, 5\}$

$$\mathcal{A}=\{\{1,2\}, \{1,3\}\}$$
$$\cup\mathcal{A}=\{1,2,3\}$$
$$f(\cup\mathcal{A})=\{1, 4, 9\}$$
$$ \bigcup \{f(A): A\in \mathcal{A}\}= \{1, 4\}\cup\{9\}$$

However:
If X and Y are a nonempty sets, $f: X\rightarrow Y$, then for any $\mathcal{A}\subseteq 2^X$
$$f(\cap\mathcal{A})\subseteq \bigcap \{f(A): A\in \mathcal{A}\}$$
E.g.
$f=x^2$\\
$X=\{-2,2\}$

$$\mathcal{A}=\{\{2\}, \{-2\}\}$$
$$\cap\mathcal{A}=\emptyset$$
$$f(\cap\mathcal{A})=\emptyset$$
$$ \bigcap \{f(A): A\in \mathcal{A}\}= \{4\}\cap\{4\}=\{4\}$$


\end{document}

